\section{Finite Rotations in 3D}
There's actually multiple ways to represent rotation in 3D space, this section explains
multiple ways of representing rotations, and specifically to rotate vectors, points and or objects in 3D eucledian space.

\subsection{Euler's Rotation Theorem}
This theorem states that given two orientations in 3D space $R_1$ and $R_2$,
then there exist an axis $\hat{A}$ and an angle $\theta$ such that it runs through a fixed point
and when applying the rotation around the axis $\hat{A}$ of $\theta$ amount to the orientation $R_1$ gives $R_2$.

In laymans' terms, you can always find a pair of axis-angle to get from one orientation to another with just 1 rotation.

\subsection{Euler-Rodrigues' Finite Rotation Formula}
Euler-Rodrigues' Finite Rotation Formula states that given a vector $\vec{v}$, we can
obtain the rotated vector $\vec{v'}$ by rotating it around an axis which is denoted by
a unit vector $\hat{k}$ of an angle of $\theta$.

\begin{center}
    $\vec{v'}=\vec{v}~cos~\theta + (\hat{k} \times \vec{v})~sin~\theta + \hat{k}(\hat{k} \cdot \vec{v})(1 - cos~\theta)$
\end{center}

where:
\begin{itemize}
    \item $\vec{v}$: is the vector to rotate.
    \item $\vec{v'}$: is the resulting rotated vector along the axis.
    \item $\hat{k}$: is the unit vector denoting the axis of rotation.
    \item $\theta$: is the angle of rotation to apply to the vector around the axis $\hat{k}$.
\end{itemize}

It is also useful to define $\hat{k}$ with respect to two non-zero vectors $\vec{a}$ and $\vec{b}$ ($\vec{a} \neq \vec{b}$) which
can denote a plane of rotation:

\begin{center}
    $\hat{k} = \frac{\vec{a} \times \vec{b}}{\|\vec{a} \times \vec{b}\|}$
\end{center}

\vspace{.5em}

\emph{
    NOTE: The formula shown above is more computationally efficient compared to
    other means of rotation.
}
